% !TeX root = ../main.tex
\section{Schanuel's Lemma and the Salvetti Complex}
\label{sec:L07-W03D2}

\begin{lemma}[Schanuel's Lemma]
    \label{lem:Schanuel}
    Suppose \(P, P'\) are projective modules, and we have exact sequences
    \[ 0 \to K \to P \xrightarrow{\pi} M \to 0 ,\] 
    \[ 0 \to K' \to P' \xrightarrow{\pi'} M \to 0 .\]
    Then  \(P\oplus K' \cong P' \oplus K\).
\end{lemma}
\begin{proof}
    Let \(Q = \{(p,p')\in P\oplus P' \mid \pi(p) = \pi'(p')\}\) be the fiber product of \(\pi, \pi'\). Define \(\a: Q \to P\) by \((p,p')\mapsto p\). Then \(\ker \a = \{(p,p') \mid p=0, \pi'(p') = \pi(p)=0\}\). But that's just \(0\oplus \ker \pi' =  K'\). Since \(\pi'\) is surjective, so \(\forall p \in P\), \(\exists p'\in P'\) such that \(\pi'(p') = \pi(p)\), that is, \(\a\) is surjective. So the following short sequence is exact:
    \[ 0 \to K' \to Q \xrightarrow{\a} P \to 0.\]

    Since \(P\) is projective, \(Q\cong P\oplus K'\) (See \cref{prop:characterization-of-projctive-module}). By the same argument, \(P'\oplus K \cong Q\).
\end{proof}

\begin{corollary}
    \label{cor:kernel-of-fg-projective-complex}
    Suppose \(P_1 \to P_0 \xrightarrow{\pi} M \to 0\) is an exact complex where \(P_0, P_1\) are finitely generated projective modules. Suppose also \(P_0'\xrightarrow{\pi'} M \to 0\) is an exact complex where \(P_0'\) is a finitely generated projective module. Then \(\ker \pi'\) is finitely generated.
\end{corollary}
\begin{proof}
    The map \(\pi\) gives a short exact sequence:
    \[0 \to \ker \pi \to P_0 \xrightarrow{\pi} M \to 0.\]
    Notice that as a submodule of a finitely generated module, \(\ker \pi\) is also finitely generated.

    We also have:
    \[0 \to \ker \pi' \to P_0' \xrightarrow{\pi'} M \to 0.\]

    By \cref{lem:Schanuel}, \(\ker \pi \oplus P_0' \cong \ker \pi' \oplus P_0\). Since \(\ker \pi, P_0', P_0\) are all finitely generated, so is \(\ker \pi'\).
\end{proof}

Recall that groups with type \(FP_n\) give rise to finitely generated projective resolutions (See \cref{def:type-FF-FP}).
\cref{prop:kernel-of-fg-projective-complex-high-dim} is a higher dimension version of \cref{lem:kernel-of-fg-projective-complex}.


\begin{proposition}
    \label{prop:kernel-of-fg-projective-complex-high-dim}
    Let \(M\) be a module. Suppose \(\{P_i\}_{i=0}^n\) is a finitely generated projective resolution of \(M\). Suppose \(\{P_i'\}_{i=0}^k\) is another finitely generated projective resolution and \(k<n\). Then \(\ker (P_k' \to P_{k-1}')\) is also finitely generated.
\end{proposition}

Recall \cref{def:type-FF-FP} and \cref{def:finite-type-KG1}.

\begin{lemma}
    \label{lem:FF1-implies-F1}
    Suppose a group \(G\) has type 
    \(FF_1\), then \(G\) is finitely generated.
\end{lemma}
\begin{proof}
    If a group \(G\)  has type   \(FF_1\) then there exists a free resolution  of \(\Z\) with \(\Z G\)-module:
    \[ F_1 \to F_0\xrightarrow{\a} \Z \to 0.\]
    We also have the map 
    \[\Z G \xrightarrow{\epsilon} \Z \to 0, \quad g\mapsto 1.\]
    By \cref{cor:kernel-of-fg-projective-complex}, \(\ker \epsilon\) is finitely generated as a \(\Z G\)-module.

    
    \begin{claim}
        Group \(G\) is finitely generated if and only if \(\ker \epsilon\) is finitely generated as a \(\Z G\)-module.
    \end{claim}
    
    Let \(S\) be a generating set of \(G\), then \(\{1-s\mid s\in S\}\) is a generating set of \(\ker \epsilon\). So \(\ker \epsilon\) is finitely generated if \(S\) is finite.
    
    Suppose now \(\ker \epsilon\) is finitely generated.
    Notice two facts:
    \begin{enumerate}
        \item \(\ker \epsilon\) is generated by \(\{1-g \mid g\in G\}\), and
        \item \(ab-1 = a(b-1)+(a-1)\).
    \end{enumerate}
    Then there exists a finite subset \(\{1-s \mid s\in S\}\) of \(\{1-g \mid g\in G\}\) that generates \(\ker \epsilon\). Take the Cayley graph \(\G\) of \(G\) with respect to \(S\), recall that \(\G\) is connected if and only if \(S\) generates \(G\). Now consider the chain complex
    \[ C_1 (\G) \xrightarrow{\partial_1} C_0 (\G) \to \Z \to 0 .\]
    Here, \(C_0(\G) =\Z G\). We verify that \(\partial_1 (e) = g-gs = g(1-s) =0\). We can show \(\ker \epsilon \subseteq \im \partial_1\) using the same strategy as in \cref{ex:compact-supported-cohomology-R}.

    It follows that \(\tilde H_0(\G)=0\), i.e. \(\G\) is connected, i.e. \(G\) is finitely generated.
\end{proof}

TODO FF2 implies FH2

TODO Salvetti complex